\section{Experimental Design}
\label{sec:experimental_design}

\subsection{Route A: parent-linked mechanism rounds}

The campaign begins from an immutable negative result: existing sparse
equation-discovery methods degraded under observation noise. Sparse
identification of nonlinear dynamics supplies the formulation
~\cite{brunton2016sindy}; PySINDy supplies a reference implementation pattern
~\cite{desilva2020pysindy}; and Operon is the genetic-programming baseline
~\cite{burlacu2020operon}. The benchmark contains ordinary and partial
differential-equation systems under several noise settings
~\cite{ziaei2025mdbench}. These citations define context only. Candidate
decisions come from the frozen local results.

An archive audit excludes systems used in two earlier panels, leaving two
mutually exclusive unseen panels of six systems. Each round uses two
development systems, one unseen panel, clean and SNR20 conditions, and three
fresh random seeds for the stochastic scientific algorithms. The candidate,
Operon baseline, and three mechanism-specific ablations produce 240 cells per
round. Outcome access occurs only after a development median improvement of at
least 0.15 and complete candidate and baseline coverage.

Round 1, \path{noise_conditioned_ensemble_sindy}, applies
noise-conditioned Savitzky--Golay derivative estimation and coefficient
bootstrap aggregation. Its ablations remove noise conditioning, aggregation,
or smoothing. Round 2, \path{spline_group_sparse_sindy}, uses analytic
smoothing-spline derivatives and cross-output group-sparse projection. Its
ablations remove spline derivatives, group projection, or adaptive
thresholding. Round 2 stores the Round 1 negative-result hash and declares a
mechanism delta before execution.

The method gate requires a system-level, failure-aware bootstrap 95\%
interval with a lower endpoint above zero, three-seed scientific
reproduction, the strong baseline, all ablations, and an idempotent rerun.
Failure on any condition writes a negative decision. Thresholds and panels
remain unchanged after outcome access.

\subsection{Route B: controlled workflow matrix}

After both Route A rounds fail, the campaign freezes a systems study. Four
sources are fresh local UCI executions~\cite{dua2019uci}; six are trace
summaries from the previously observed equation-discovery campaign. Each
source record contains hashes, validated scalar metrics, a truth label, and a
use restriction. UCI runs show ingestion of newly computed evidence. MDBench
traces connect the state-machine study to the failed rounds but cannot serve
as new method holdouts.

Table~\ref{tab:task_matrix} lists the ten tasks. One frozen workflow state is
injected into each otherwise valid source task. Faults cover source integrity,
metric direction, configuration freeze, missing evidence maps, unsupported
claims, missing repetition coverage, failed cells, and claim-evidence
mismatch. Three tasks require a prior \vault{} record. One task in each family
has no injected fault.

\begin{table*}[t]
\centering
\caption{Frozen tasks and controlled workflow states. Source evidence is
observed data, while the workflow fault is deliberately introduced. MDBench
rows are behavior tests rather than method holdouts.}
\label{tab:task_matrix}
\small
\begin{tabular}{llll}
\toprule
Family & Source task & Initial workflow state & Required repair \\
\midrule
UCI & PenDigits prototypes & none & retain valid evidence \\
UCI & Letter prototypes & stale source hash & rehash source \\
UCI & Spambase prototypes & wrong metric direction & bind direction \\
UCI & Skin prototypes & missing evidence map & restore from \vault{} \\
\midrule
MDBench & binocular rivalry & none & retain valid evidence \\
MDBench & interacting magnets & unfrozen configuration & freeze configuration \\
MDBench & oscillator death & unsupported claim & evidence-bound decision \\
MDBench & population growth & incomplete repetitions & resume checkpoint \\
MDBench & RC circuit & failed experiment cell & classify and recover \\
MDBench & Van der Pol & claim-evidence mismatch & rewrite claim \\
\bottomrule
\end{tabular}
\end{table*}

\subsection{Controllers and component removals}

The \oneshot{} controller emits a result from the initial state. The
\executeonce{} controller plans once and repairs only a fault visible during
planning. The \systemloop{} freezes a contract, executes, classifies an
observed failure, retrieves permitted \vault{} context, applies the task's
registered repair, reruns once, and checks claim support. These policies are
deterministic in Route B. A local model supplies a structured policy summary
but cannot change a transition or score.

Four registered component removals isolate the contracts. Removing \vault{}
memory blocks repairs that need a prior evidence card. Removing failure
feedback prevents the second action after a runtime diagnosis. Removing
preregistration makes the completion contract false. Removing the evidence
gate hides claim-support failures and permits an unsupported report claim.

Seven modes run on \TaskCount{} tasks with \IdempotencyRepeatCount{}
deterministic repetition identifiers. This produces \MainCellCount{} main
cells and \AblationCellCount{} component-removal cells, or \CellCount{} cells
in total. Each repetition recomputes the scientific hash. The transition
policy does not sample stochastic repair behavior, so equal outcomes across
the three identifiers test idempotency and accidental identifier dependence.
They do not represent three independent controller trajectories.

\subsection{Endpoints and additive task-level analysis}

Task success requires valid source identity, completed execution, a complete
evidence map, reproduction readiness, supported report claims, and, for the
full loop, a frozen contract. Recovery is a cell-level engineering measure:
an initially failed cell takes a second action and then satisfies the task
contract. An unsupported claim has a polarity or scope inconsistent with its
source record. Intervention counts any human research decision after campaign
start. External cost counts paid API and cloud-compute use.

The primary independent unit for the additive analysis is task. After
confirming that the three deterministic repetitions agree within each task,
let \(F_t\) and \(E_t\) denote full-loop and execute-once task success. We
define
\[
d_t=F_t-E_t.
\]
The ten values are sampled with replacement for \BootstrapResamples{}
task-bootstrap resamples using the registered random seed. We report the
percentile interval and an exact sign test over nonzero task differences;
ties are retained in the task inventory and omitted only from the conditional
sign count. The bootstrap follows the nonparametric procedure
~\cite{efron1979bootstrap}, and the task-level unit follows cross-data-set
comparison practice~\cite{demsar2006}.

The original internal gate averaged 30 task-repetition differences and
reported [\HistoricalCellCILower{}, \HistoricalCellCIUpper{}]. That interval
is retained as historical protocol evidence. Because repeated deterministic
outcomes do not create new independent tasks, it is not used for a
publication-facing effect statement. The task-level analysis was added after
the field-currency audit and does not replace the original preregistration.
